\documentclass[11pt,a4paper]{article}
\usepackage[utf8]{inputenc}
\usepackage[T1]{fontenc}
\usepackage[spanish,es-minimal,es-noshorthands]{babel}
\AtBeginDocument{\renewcommand{\sin}{\operatorname{sen}}\renewcommand{\arcsin}{\operatorname{arcsen}}}
\usepackage{lmodern}
\usepackage[sfdefault,lining]{FiraSans}
\usepackage[a4paper,top=2.3cm,bottom=2.2cm,left=1.8cm,right=1.8cm,headheight=1.3cm,headsep=0.4cm,footskip=1cm]{geometry}
\usepackage{amsmath,amssymb}
\usepackage{icomma}
\usepackage{graphicx}
\usepackage[table]{xcolor}
\usepackage{tikz}
\usepackage{fontawesome5}
\usepackage{ccicons}
\usepackage{fancyhdr}
\usepackage{lastpage}
\usepackage{tasks}
\usepackage{needspace}
\usepackage[most]{tcolorbox}
\usepackage{array}
\usepackage[hidelinks]{hyperref}
\AtBeginDocument{\hypersetup{pdftitle={Soluciones de la ficha de la unidad 4: Ecuaciones de 1.er y 2.º grado. 2º ESO},pdfsubject={Matemáticas 2º ESO},pdfkeywords={ecuaciones, identidades, ecuaciones de primer grado, paréntesis, denominadores, problemas, ecuaciones de segundo grado, incompletas, fórmula, discriminante}}}

\definecolor{navy}{HTML}{1B2748}
\definecolor{accent}{HTML}{E8603A}
\definecolor{teal}{HTML}{1F9E90}
\definecolor{soft}{HTML}{EEF1F8}
\definecolor{muted}{HTML}{5D6685}

\setlength{\parindent}{0pt}
\setlength{\parskip}{3pt}
\settasks{label=\alph*), label-format=\bfseries\color{navy}, label-width=1.3em, item-indent=1.8em, column-sep=0.6em, after-item-skip=5pt, before-skip=3pt, after-skip=4pt}

% ---------- capçalera i peu ----------
\newcommand{\llicencia}{\href{https://creativecommons.org/licenses/by-nc-sa/4.0/deed.es}{\ccbyncsa}\ \ Licencia Creative Commons Atribución-NoComercial-CompartirIgual 4.0 Internacional (CC BY-NC-SA 4.0)}
\pagestyle{fancy}
\fancyhf{}
\fancyhead[L]{\raisebox{-0.35\height}{\includegraphics[height=0.95cm]{logo-gva.png}}}
\fancyhead[C]{\small\color{muted}Soluciones de la ficha · Unidad 4: Ecuaciones de 1.\textsuperscript{er} y 2.\textsuperscript{o} grado · 2\textsuperscript{o} ESO}
\fancyhead[R]{\raisebox{-0.35\height}{\includegraphics[height=0.95cm]{logo-vermellar.png}}}
\fancyfoot[L]{\scriptsize\color{muted}\llicencia}
\fancyfoot[R]{\small\color{muted}\thepage}
\renewcommand{\headrulewidth}{0.4pt}
\renewcommand{\headrule}{\color{muted!40}\hrule height \headrulewidth}
\fancypagestyle{primera}{\fancyhf{}\renewcommand{\headrulewidth}{0pt}
  \fancyfoot[L]{\scriptsize\color{muted}\llicencia}
  \fancyfoot[R]{\small\color{muted}\thepage}}

% ---------- nivells i exercicis ----------
\newcommand{\estrelles}[1]{{\color{accent}\ifcase#1\or\faStar\faStar[regular]\faStar[regular]\or\faStar\faStar\faStar[regular]\or\faStar\faStar\faStar\fi}}
\newcounter{exer}
\newcommand{\ex}[1]{\par\needspace{4\baselineskip}\medskip\stepcounter{exer}\noindent{\color{navy}\bfseries\theexer.}\ \ {\small\estrelles{#1}}\ \ }

\newcommand{\seccio}[2]{\par\needspace{10\baselineskip}\bigskip
  \begin{tcolorbox}[colback=navy,colframe=navy,coltext=white,arc=2mm,boxrule=0pt,left=3mm,right=3mm,top=1.2mm,bottom=1.2mm]
  {\large\bfseries #1\quad #2}
  \end{tcolorbox}\vspace{-2pt}}

\newcolumntype{C}[1]{>{\centering\arraybackslash}p{#1}}
% quadre de recordatori (només a la fitxa d'exercicis)
\newcommand{\recorda}[1]{\par\smallskip\begin{tcolorbox}[colback=accent!7,colframe=accent!55,arc=2mm,boxrule=0.5pt,left=2.5mm,right=2.5mm,top=1mm,bottom=1mm]\small{\color{accent}\faLightbulb[regular]}\ \textbf{Recuerda:}\ #1\end{tcolorbox}}
\newcommand{\res}[1]{\boxed{#1}}\newcommand{\rest}[1]{\fbox{#1}}

\def\fmneed{12}
% =====================================================================
% Capçalera FentMates de les fitxes: un color per a cada curs.
% Els colors són els de la pàgina inicial (LEVELS a assets/js/materials.js);
% el logo ve de _eines/logo/pdf/ (python3 _eines/logo/genera.py).
%
% Ús, al final del preàmbul de la unitat:
%   % =====================================================================
% Capçalera FentMates de les fitxes: un color per a cada curs.
% Els colors són els de la pàgina inicial (LEVELS a assets/js/materials.js);
% el logo ve de _eines/logo/pdf/ (python3 _eines/logo/genera.py).
%
% Ús, al final del preàmbul de la unitat:
%   % =====================================================================
% Capçalera FentMates de les fitxes: un color per a cada curs.
% Els colors són els de la pàgina inicial (LEVELS a assets/js/materials.js);
% el logo ve de _eines/logo/pdf/ (python3 _eines/logo/genera.py).
%
% Ús, al final del preàmbul de la unitat:
%   \input{../comu/fentmates.tex}
%   \fmcurs{1eso}
% i, al començament del document:
%   \capcalera{Matemàtiques}{1r ESO}{Fitxa d'exercicis · Unitat 2}{Divisibilitat}
% Opcional abans de \input: \def\fmneed{14} (espai mínim abans de cada secció).
% =====================================================================
\makeatletter
\def\fm@logodir{../../logo/pdf}
\providecommand{\fmneed}{10}
\def\fm@cursnom{}

% clau del curs, color i color del text damunt del color
\newcommand{\fm@curs}[3]{\@namedef{fm@color@#1}{#2}\@namedef{fm@text@#1}{#3}}
\fm@curs{1eso}{E6B84A}{navy}
\fm@curs{2eso}{3B82F6}{white}
\fm@curs{3eso}{FF7A50}{white}
\fm@curs{4eso-a}{A78BFA}{white}
\fm@curs{4eso-b}{6D28D9}{white}
\fm@curs{1batx-mi}{9B1C1C}{white}
\fm@curs{1batx-ccss}{EF6B6B}{white}
\fm@curs{2batx-mii}{1E6B47}{white}
\fm@curs{2batx-ccss}{2BB3A3}{white}

\newcommand{\fmcurs}[1]{%
  \@ifundefined{fm@color@#1}{\PackageError{fentmates}{Curs desconegut: #1}{}}{}%
  \edef\fm@tmp{\noexpand\definecolor{curs}{HTML}{\@nameuse{fm@color@#1}}%
               \noexpand\colorlet{curstext}{\@nameuse{fm@text@#1}}}\fm@tmp
  \colorlet{cursfosc}{curs!72!black}%
  \def\fm@logo{\fm@logodir/vertical-#1.pdf}%
  \def\fm@marca{\fm@logodir/marca-#1.pdf}}

% ---------- capçalera de la primera pàgina ----------
% #1 matèria (petit), #2 curs, #3 tipus de document i unitat, #4 títol
\newcommand{\capcalera}[4]{%
  \gdef\fm@cursnom{#2}%
  \par\noindent
  \begin{tcolorbox}[colback=curs,colframe=curs,boxrule=0pt,arc=2.2mm,width=3.5cm,height=2.15cm,
      valign=center,left=3mm,right=2mm,top=1mm,bottom=1mm,nobeforeafter,box align=center]
    \color{curstext}{\footnotesize #1}\par\vspace{1pt}{\fontsize{19}{21}\selectfont\bfseries #2\par}%
  \end{tcolorbox}\hspace{5mm}%
  \begin{minipage}[c]{\dimexpr\linewidth-3.5cm-5mm-2.9cm\relax}
    {\color{cursfosc}\small\bfseries\MakeUppercase{#3}}\par\vspace{2pt}
    {\color{navy}\fontsize{20}{23}\selectfont\bfseries #4\par}
  \end{minipage}\hfill
  \begin{minipage}[c]{2.6cm}\raggedleft\includegraphics[width=2.6cm]{\fm@logo}\end{minipage}\par
  \vspace{2pt}}

% ---------- barres de secció del color del curs ----------
\renewcommand{\seccio}[2]{\par\needspace{\fmneed\baselineskip}\bigskip
  \begin{tcolorbox}[colback=curs,colframe=curs,coltext=curstext,arc=2mm,boxrule=0pt,left=3mm,right=3mm,top=1.2mm,bottom=1.2mm]
  {\large\bfseries #1\quad #2}
  \end{tcolorbox}\vspace{-2pt}}

% ---------- estrelles de dificultat del color del curs (a Batxillerat no se'n posen) ----------
\@ifundefined{estrelles}{}{%
  \renewcommand{\estrelles}[1]{{\color{curs}\ifcase#1\or\faStar\faStar[regular]\faStar[regular]\or\faStar\faStar\faStar[regular]\or\faStar\faStar\faStar\fi}}}

% ---------- capçalera i peu de les pàgines ----------
% primera pàgina: logos institucionals al peu
\fancypagestyle{primera}{\fancyhf{}\renewcommand{\headrulewidth}{0pt}
  \fancyfoot[L]{\raisebox{-0.32\height}{\includegraphics[height=0.78cm]{logo-gva.png}}\hspace{5mm}%
    \raisebox{-0.32\height}{\includegraphics[height=0.78cm]{logo-vermellar.png}}}
  \fancyfoot[R]{\scriptsize\color{muted}\href{https://creativecommons.org/licenses/by-nc-sa/4.0/deed.ca}{\ccbyncsa}\ \ CC BY-NC-SA 4.0 · fentmates.org\qquad\small\thepage}}
% la resta: el símbol a l'esquerra i el curs a la dreta, amb el títol al centre com abans
\fancyhead[L]{\raisebox{-0.3\height}{\includegraphics[height=0.62cm]{\fm@marca}}}
\fancyhead[R]{\small\bfseries\color{cursfosc}\fm@cursnom}
\renewcommand{\headrule}{\color{curs!45}\hrule height \headrulewidth}
\makeatother

%   \fmcurs{1eso}
% i, al començament del document:
%   \capcalera{Matemàtiques}{1r ESO}{Fitxa d'exercicis · Unitat 2}{Divisibilitat}
% Opcional abans de \input: \def\fmneed{14} (espai mínim abans de cada secció).
% =====================================================================
\makeatletter
\def\fm@logodir{../../logo/pdf}
\providecommand{\fmneed}{10}
\def\fm@cursnom{}

% clau del curs, color i color del text damunt del color
\newcommand{\fm@curs}[3]{\@namedef{fm@color@#1}{#2}\@namedef{fm@text@#1}{#3}}
\fm@curs{1eso}{E6B84A}{navy}
\fm@curs{2eso}{3B82F6}{white}
\fm@curs{3eso}{FF7A50}{white}
\fm@curs{4eso-a}{A78BFA}{white}
\fm@curs{4eso-b}{6D28D9}{white}
\fm@curs{1batx-mi}{9B1C1C}{white}
\fm@curs{1batx-ccss}{EF6B6B}{white}
\fm@curs{2batx-mii}{1E6B47}{white}
\fm@curs{2batx-ccss}{2BB3A3}{white}

\newcommand{\fmcurs}[1]{%
  \@ifundefined{fm@color@#1}{\PackageError{fentmates}{Curs desconegut: #1}{}}{}%
  \edef\fm@tmp{\noexpand\definecolor{curs}{HTML}{\@nameuse{fm@color@#1}}%
               \noexpand\colorlet{curstext}{\@nameuse{fm@text@#1}}}\fm@tmp
  \colorlet{cursfosc}{curs!72!black}%
  \def\fm@logo{\fm@logodir/vertical-#1.pdf}%
  \def\fm@marca{\fm@logodir/marca-#1.pdf}}

% ---------- capçalera de la primera pàgina ----------
% #1 matèria (petit), #2 curs, #3 tipus de document i unitat, #4 títol
\newcommand{\capcalera}[4]{%
  \gdef\fm@cursnom{#2}%
  \par\noindent
  \begin{tcolorbox}[colback=curs,colframe=curs,boxrule=0pt,arc=2.2mm,width=3.5cm,height=2.15cm,
      valign=center,left=3mm,right=2mm,top=1mm,bottom=1mm,nobeforeafter,box align=center]
    \color{curstext}{\footnotesize #1}\par\vspace{1pt}{\fontsize{19}{21}\selectfont\bfseries #2\par}%
  \end{tcolorbox}\hspace{5mm}%
  \begin{minipage}[c]{\dimexpr\linewidth-3.5cm-5mm-2.9cm\relax}
    {\color{cursfosc}\small\bfseries\MakeUppercase{#3}}\par\vspace{2pt}
    {\color{navy}\fontsize{20}{23}\selectfont\bfseries #4\par}
  \end{minipage}\hfill
  \begin{minipage}[c]{2.6cm}\raggedleft\includegraphics[width=2.6cm]{\fm@logo}\end{minipage}\par
  \vspace{2pt}}

% ---------- barres de secció del color del curs ----------
\renewcommand{\seccio}[2]{\par\needspace{\fmneed\baselineskip}\bigskip
  \begin{tcolorbox}[colback=curs,colframe=curs,coltext=curstext,arc=2mm,boxrule=0pt,left=3mm,right=3mm,top=1.2mm,bottom=1.2mm]
  {\large\bfseries #1\quad #2}
  \end{tcolorbox}\vspace{-2pt}}

% ---------- estrelles de dificultat del color del curs (a Batxillerat no se'n posen) ----------
\@ifundefined{estrelles}{}{%
  \renewcommand{\estrelles}[1]{{\color{curs}\ifcase#1\or\faStar\faStar[regular]\faStar[regular]\or\faStar\faStar\faStar[regular]\or\faStar\faStar\faStar\fi}}}

% ---------- capçalera i peu de les pàgines ----------
% primera pàgina: logos institucionals al peu
\fancypagestyle{primera}{\fancyhf{}\renewcommand{\headrulewidth}{0pt}
  \fancyfoot[L]{\raisebox{-0.32\height}{\includegraphics[height=0.78cm]{logo-gva.png}}\hspace{5mm}%
    \raisebox{-0.32\height}{\includegraphics[height=0.78cm]{logo-vermellar.png}}}
  \fancyfoot[R]{\scriptsize\color{muted}\href{https://creativecommons.org/licenses/by-nc-sa/4.0/deed.ca}{\ccbyncsa}\ \ CC BY-NC-SA 4.0 · fentmates.org\qquad\small\thepage}}
% la resta: el símbol a l'esquerra i el curs a la dreta, amb el títol al centre com abans
\fancyhead[L]{\raisebox{-0.3\height}{\includegraphics[height=0.62cm]{\fm@marca}}}
\fancyhead[R]{\small\bfseries\color{cursfosc}\fm@cursnom}
\renewcommand{\headrule}{\color{curs!45}\hrule height \headrulewidth}
\makeatother

%   \fmcurs{1eso}
% i, al començament del document:
%   \capcalera{Matemàtiques}{1r ESO}{Fitxa d'exercicis · Unitat 2}{Divisibilitat}
% Opcional abans de \input: \def\fmneed{14} (espai mínim abans de cada secció).
% =====================================================================
\makeatletter
\def\fm@logodir{../../logo/pdf}
\providecommand{\fmneed}{10}
\def\fm@cursnom{}

% clau del curs, color i color del text damunt del color
\newcommand{\fm@curs}[3]{\@namedef{fm@color@#1}{#2}\@namedef{fm@text@#1}{#3}}
\fm@curs{1eso}{E6B84A}{navy}
\fm@curs{2eso}{3B82F6}{white}
\fm@curs{3eso}{FF7A50}{white}
\fm@curs{4eso-a}{A78BFA}{white}
\fm@curs{4eso-b}{6D28D9}{white}
\fm@curs{1batx-mi}{9B1C1C}{white}
\fm@curs{1batx-ccss}{EF6B6B}{white}
\fm@curs{2batx-mii}{1E6B47}{white}
\fm@curs{2batx-ccss}{2BB3A3}{white}

\newcommand{\fmcurs}[1]{%
  \@ifundefined{fm@color@#1}{\PackageError{fentmates}{Curs desconegut: #1}{}}{}%
  \edef\fm@tmp{\noexpand\definecolor{curs}{HTML}{\@nameuse{fm@color@#1}}%
               \noexpand\colorlet{curstext}{\@nameuse{fm@text@#1}}}\fm@tmp
  \colorlet{cursfosc}{curs!72!black}%
  \def\fm@logo{\fm@logodir/vertical-#1.pdf}%
  \def\fm@marca{\fm@logodir/marca-#1.pdf}}

% ---------- capçalera de la primera pàgina ----------
% #1 matèria (petit), #2 curs, #3 tipus de document i unitat, #4 títol
\newcommand{\capcalera}[4]{%
  \gdef\fm@cursnom{#2}%
  \par\noindent
  \begin{tcolorbox}[colback=curs,colframe=curs,boxrule=0pt,arc=2.2mm,width=3.5cm,height=2.15cm,
      valign=center,left=3mm,right=2mm,top=1mm,bottom=1mm,nobeforeafter,box align=center]
    \color{curstext}{\footnotesize #1}\par\vspace{1pt}{\fontsize{19}{21}\selectfont\bfseries #2\par}%
  \end{tcolorbox}\hspace{5mm}%
  \begin{minipage}[c]{\dimexpr\linewidth-3.5cm-5mm-2.9cm\relax}
    {\color{cursfosc}\small\bfseries\MakeUppercase{#3}}\par\vspace{2pt}
    {\color{navy}\fontsize{20}{23}\selectfont\bfseries #4\par}
  \end{minipage}\hfill
  \begin{minipage}[c]{2.6cm}\raggedleft\includegraphics[width=2.6cm]{\fm@logo}\end{minipage}\par
  \vspace{2pt}}

% ---------- barres de secció del color del curs ----------
\renewcommand{\seccio}[2]{\par\needspace{\fmneed\baselineskip}\bigskip
  \begin{tcolorbox}[colback=curs,colframe=curs,coltext=curstext,arc=2mm,boxrule=0pt,left=3mm,right=3mm,top=1.2mm,bottom=1.2mm]
  {\large\bfseries #1\quad #2}
  \end{tcolorbox}\vspace{-2pt}}

% ---------- estrelles de dificultat del color del curs (a Batxillerat no se'n posen) ----------
\@ifundefined{estrelles}{}{%
  \renewcommand{\estrelles}[1]{{\color{curs}\ifcase#1\or\faStar\faStar[regular]\faStar[regular]\or\faStar\faStar\faStar[regular]\or\faStar\faStar\faStar\fi}}}

% ---------- capçalera i peu de les pàgines ----------
% primera pàgina: logos institucionals al peu
\fancypagestyle{primera}{\fancyhf{}\renewcommand{\headrulewidth}{0pt}
  \fancyfoot[L]{\raisebox{-0.32\height}{\includegraphics[height=0.78cm]{logo-gva.png}}\hspace{5mm}%
    \raisebox{-0.32\height}{\includegraphics[height=0.78cm]{logo-vermellar.png}}}
  \fancyfoot[R]{\scriptsize\color{muted}\href{https://creativecommons.org/licenses/by-nc-sa/4.0/deed.ca}{\ccbyncsa}\ \ CC BY-NC-SA 4.0 · fentmates.org\qquad\small\thepage}}
% la resta: el símbol a l'esquerra i el curs a la dreta, amb el títol al centre com abans
\fancyhead[L]{\raisebox{-0.3\height}{\includegraphics[height=0.62cm]{\fm@marca}}}
\fancyhead[R]{\small\bfseries\color{cursfosc}\fm@cursnom}
\renewcommand{\headrule}{\color{curs!45}\hrule height \headrulewidth}
\makeatother

\fmcurs{2eso}

\begin{document}
\thispagestyle{primera}

\capcalera{Matemáticas}{2\textsuperscript{o} ESO}{Soluciones de la ficha · Unidad 4}{Ecuaciones de 1.\textsuperscript{er} y 2.\textsuperscript{o} grado}
\vspace{8pt}

\begin{tcolorbox}[colback=soft,colframe=soft,arc=2mm,left=3mm,right=3mm,top=1.5mm,bottom=1.5mm]
\small Intenta resolver cada ejercicio antes de mirar su solución. Si no te sale el mismo resultado, repasa los pasos o consulta los apuntes interactivos de la unidad: encontrarás calculadoras que resuelven paso a paso cualquier ecuación de primer o de segundo grado.
\end{tcolorbox}

% =====================================================
\seccio{1}{Ecuaciones: elementos y soluciones}

\ex{1} \begin{tasks}(2)
  \task \rest{Identidad}
  \task \rest{Ecuación} ($x=2$)
  \task \rest{Identidad}
  \task \rest{Ecuación} ($x=\pm1$)
  \task \rest{Identidad}
  \task \rest{Ecuación} ($x=1$)
  \task \rest{Identidad}
  \task \rest{Ecuación} sin solución ($-2\neq-1$)
\end{tasks}

\ex{1} \begin{tasks}(2)
  \task $4=4$: \rest{sí}
  \task $5=5$: \rest{sí}
  \task $8\neq9$: \rest{no}
  \task $5=5$: \rest{sí}
  \task $6=6$: \rest{sí}
  \task $4\neq7$: \rest{no}
\end{tasks}

\ex{1} 1.\textsuperscript{er} miembro: $4x^2+x-3$; 2.\textsuperscript{o} miembro: $2x+5$; términos: $4x^2,\ x,\ -3,\ 2x,\ 5$; incógnita: $x$; \rest{grado 2}.

\ex{2} \begin{tasks}(2)
  \task $4x=-8$: \rest{grado 1}
  \task $-4x=0$: \rest{grado 1}
  \task $-x^2+3x+1=0$: \rest{grado 2}
  \task $x^2-3x-1=0$: \rest{grado 2}
\end{tasks}

\ex{2} Por ejemplo, \rest{$x+2=0$} \ (o $3x+1=-5$\ldots) \ y \ \rest{$x^2-9=0$} \ (o $x^2=9$\ldots).

\ex{3} \begin{tasks}(1)
  \task $3k-4=6+k\Rightarrow 2k=10\Rightarrow\res{k=5}$
  \task $1-k=k+5\Rightarrow -2k=4\Rightarrow\res{k=-2}$
\end{tasks}


% =====================================================
\seccio{2}{Ecuaciones de primer grado}

\ex{1} \begin{tasks}(2)
  \task $\res{x=5}$
  \task $\res{x=2}$
  \task $\res{x=-7}$
  \task $\res{x=-5}$
  \task $\text{m.c.m.}=5:\ \ \res{x=15}$
  \task $\res{x=-6}$
  \task $\res{x=\tfrac{7}{2}}$
  \task $\text{m.c.m.}=4:\ \ \res{x=-8}$
\end{tasks}

\ex{1} \begin{tasks}(3)
  \task $5x=15\ \Rightarrow\ \res{x=3}$
  \task $-2x=-8\ \Rightarrow\ \res{x=4}$
  \task $3x=-6\ \Rightarrow\ \res{x=-2}$
  \task $-x=3\ \Rightarrow\ \res{x=-3}$
  \task $4x=-12\ \Rightarrow\ \res{x=-3}$
  \task $5x=-10\ \Rightarrow\ \res{x=-2}$
\end{tasks}

\ex{1} \begin{tasks}(2)
  \task $2x=8\ \Rightarrow\ \res{x=4}$
  \task $-4x=-20\ \Rightarrow\ \res{x=5}$
  \task $-3x=12\ \Rightarrow\ \res{x=-4}$
  \task $-3x=0\ \Rightarrow\ \res{x=0}$
  \task $5x=-15\ \Rightarrow\ \res{x=-3}$
  \task $-5x=-10\ \Rightarrow\ \res{x=2}$
\end{tasks}

\ex{2} \begin{tasks}(2)
  \task $3x=12\ \Rightarrow\ \res{x=4}$
  \task $-3x=-12\ \Rightarrow\ \res{x=4}$
  \task $-4x=-16\ \Rightarrow\ \res{x=4}$
  \task $-x=6\ \Rightarrow\ \res{x=-6}$
\end{tasks}

\ex{2} \begin{tasks}(2)
  \task $4x=6\ \Rightarrow\ \res{x=\tfrac{3}{2}}$
  \task $3x=5\ \Rightarrow\ \res{x=\tfrac{5}{3}}$
  \task $-8x=-4\ \Rightarrow\ \res{x=\tfrac{1}{2}}$
  \task $6x=-10\ \Rightarrow\ \res{x=-\tfrac{5}{3}}$
\end{tasks}

\ex{3} \begin{tasks}(1)
  \task $6x=18\ \Rightarrow\ \res{x=3}$
  \task $-5x=-20\ \Rightarrow\ \res{x=4}$
  \task $5x=8\ \Rightarrow\ \res{x=\tfrac{8}{5}}$
  \task $-3x=-14\ \Rightarrow\ \res{x=\tfrac{14}{3}}$
\end{tasks}

\ex{3} El $2x$ suma, y tiene que pasar \textbf{restando}; el $-3$ resta, y tiene que pasar \textbf{sumando}. \ $5x-2x=9+3\Rightarrow 3x=12\Rightarrow\res{x=4}$. \ Comprobación: $5\cdot4-3=17=2\cdot4+9$.


% =====================================================
\seccio{3}{Ecuaciones con paréntesis}

\ex{1} \begin{tasks}(2)
  \task $3x-6=12\ \Rightarrow\ 3x=18\ \Rightarrow\ \res{x=6}$
  \task $5x+5=2x-7\ \Rightarrow\ 3x=-12\ \Rightarrow\ \res{x=-4}$
  \task $2x-8=x+1\ \Rightarrow\ \res{x=9}$
  \task $8x-4=3x+6\ \Rightarrow\ 5x=10\ \Rightarrow\ \res{x=2}$
\end{tasks}

\ex{1} \begin{tasks}(2)
  \task $-x+1=2x-5\ \Rightarrow\ -3x=-6\ \Rightarrow\ \res{x=2}$
  \task $2x-5=3\ \Rightarrow\ 2x=8\ \Rightarrow\ \res{x=4}$
  \task $-2x+8=x+2\ \Rightarrow\ -3x=-6\ \Rightarrow\ \res{x=2}$
  \task $-x+6=2x\ \Rightarrow\ -3x=-6\ \Rightarrow\ \res{x=2}$
\end{tasks}

\ex{2} \begin{tasks}(1)
  \task $4x-11=5x-13\ \Rightarrow\ -x=-2\ \Rightarrow\ \res{x=2}$
  \task $2x-1=4x-7\ \Rightarrow\ -2x=-6\ \Rightarrow\ \res{x=3}$
  \task $3x-7=2x+2\ \Rightarrow\ \res{x=9}$
  \task $-x+9=4x-6\ \Rightarrow\ -5x=-15\ \Rightarrow\ \res{x=3}$
\end{tasks}

\ex{2} \begin{tasks}(1)
  \task $2x+6=2x+6\ \Rightarrow\ 0x=0\ \Rightarrow\ \res{\text{infinitas soluciones}}$
  \task $3x+2=3x+4\ \Rightarrow\ 0x=2\ \Rightarrow\ \res{\text{sin solución}}$
  \task $-2=-2\ \Rightarrow\ 0x=0\ \Rightarrow\ \res{\text{infinitas soluciones}}$
  \task $x+3=2x+1\ \Rightarrow\ -x=-2\ \Rightarrow\ \res{x=2}$
\end{tasks}

\ex{2} \begin{tasks}(1)
  \task $4x=12\ \Rightarrow\ \res{x=3}$
  \task $3x=2x+5\ \Rightarrow\ \res{x=5}$
  \task $x^{2}-4=x^{2}+x-7\ \Rightarrow\ -x=-3\ \Rightarrow\ \res{x=3}$
\end{tasks}

\ex{3} \begin{tasks}(1)
  \task $-2x+3=7\ \Rightarrow\ -2x=4\ \Rightarrow\ \res{x=-2}$
  \task $-4x+2=-7x+8\ \Rightarrow\ 3x=6\ \Rightarrow\ \res{x=2}$
\end{tasks}

\ex{3} $2(3+a)=12\Rightarrow 6+2a=12\Rightarrow\res{a=3}$


% =====================================================
\seccio{4}{Ecuaciones con denominadores}

\ex{1} \begin{tasks}(2)
  \task $\text{m.c.m.}=3:\ \ \res{x=12}$
  \task $\text{m.c.m.}=4:\ \ 2x+x=24\ \Rightarrow\ 3x=24\ \Rightarrow\ \res{x=8}$
  \task $\text{m.c.m.}=3:\ \ 3x-x=30\ \Rightarrow\ 2x=30\ \Rightarrow\ \res{x=15}$
  \task $\text{m.c.m.}=5:\ \ 2x-5=15\ \Rightarrow\ 2x=20\ \Rightarrow\ \res{x=10}$
\end{tasks}

\ex{1} \begin{tasks}(2)
  \task $\text{m.c.m.}=3:\ \ (x+2)=12\ \Rightarrow\ x+2=12\ \Rightarrow\ \res{x=10}$
  \task $\text{m.c.m.}=4:\ \ 2(x-1)=(x+3)\ \Rightarrow\ 2x-2=x+3\ \Rightarrow\ \res{x=5}$
  \task $\text{m.c.m.}=4:\ \ (3x-1)=8\ \Rightarrow\ 3x-1=8\ \Rightarrow\ 3x=9\ \Rightarrow\ \res{x=3}$
  \task $\text{m.c.m.}=12:\ \ 2x+12=3x\ \Rightarrow\ -x=-12\ \Rightarrow\ \res{x=12}$
\end{tasks}

\ex{2} \begin{tasks}(1)
  \task $\text{m.c.m.}=6:\ \ 3(x+1)-2(x-1)=12\ \Rightarrow\ x+5=12\ \Rightarrow\ \res{x=7}$
  \task $\text{m.c.m.}=12:\ \ 3(x-2)-2(x-3)=6\ \Rightarrow\ \res{x=6}$
  \task $\text{m.c.m.}=6:\ \ 2(2x+1)-3(x-2)=18\ \Rightarrow\ x+8=18\ \Rightarrow\ \res{x=10}$
  \task $\text{m.c.m.}=10:\ \ 10x-2(x+2)=5(x+6)\ \Rightarrow\ 8x-4=5x+30\ \Rightarrow\ 3x=34\ \Rightarrow\ \res{x=\tfrac{34}{3}}$
\end{tasks}

\ex{2} \begin{tasks}(1)
  \task $\text{m.c.m.}=6:\ \ 2(x-1)+(x+2)=6x-12\ \Rightarrow\ 3x=6x-12\ \Rightarrow\ -3x=-12\ \Rightarrow\ \res{x=4}$
  \task $\text{m.c.m.}=4:\ \ 12-(x-1)=2(x+5)\ \Rightarrow\ -x+13=2x+10\ \Rightarrow\ -3x=-3\ \Rightarrow\ \res{x=1}$
\end{tasks}

\ex{3} \begin{tasks}(1)
  \task $\text{m.c.m.}=12:\ \ 9(x-2)-8(x+1)=2(x-10)\ \Rightarrow\ x-26=2x-20\ \Rightarrow\ -x=6\ \Rightarrow\ \res{x=-6}$
  \task $\text{m.c.m.}=12:\ \ 6x-4(x-1)-3(x+1)=(1-x)\ \Rightarrow\ -x+1=-x+1\ \Rightarrow\ 0x=0\ \Rightarrow\ \res{\text{infinitas soluciones}}$
  \task $\text{m.c.m.}=4:\ \ 2(x+1)-(x+3)=x-4\ \Rightarrow\ x-1=x-4\ \Rightarrow\ 0x=-3\ \Rightarrow\ \res{\text{sin solución}}$
\end{tasks}

\ex{3} El $-$ delante de la segunda fracción afecta a todo el numerador: $-3(x+1)=-3x-3$. \ $2x-4-3x-3=6\Rightarrow -x=13\Rightarrow\res{x=-13}$.


% =====================================================
\seccio{5}{Problemas con ecuaciones de primer grado}

\ex{1} $x+(x+1)=75\Rightarrow 2x=74\Rightarrow x=37$: \ \rest{37 y 38}

\ex{1} $x+3x=56\Rightarrow 4x=56\Rightarrow\res{x=14}$

\ex{1} $40+x=2(12+x)\Rightarrow 40+x=24+2x\Rightarrow\res{x=16\text{ años}}$

\ex{2} $2x+2(2x+3)=54\Rightarrow 6x+6=54\Rightarrow x=8$: \ \rest{8 m y 19 m}

\ex{2} Gallinas $x$: \ $2x+4(40-x)=110\Rightarrow -2x=-50\Rightarrow x=25$: \ \rest{25 gallinas y 15 conejos}

\ex{2} Monedas de 2 €: $x$. \ $2x+(30-x)=46\Rightarrow\res{x=16}$: \ \rest{16 de 2 € y 14 de 1 €}

\ex{2} $3x+12=2(x+12)\Rightarrow\res{x=12}$: \ hijo \rest{12 años}, madre \rest{36 años}

\ex{2} $0{,}85x=68\Rightarrow x=\tfrac{68}{0{,}85}=\res{80\text{ €}}$

\ex{3} $x-\tfrac{x}{3}-\tfrac{x}{4}=25\xrightarrow{\cdot12}12x-4x-3x=300\Rightarrow 5x=300\Rightarrow\res{60\text{ €}}$

\ex{3} $x+2x+(2x+30)=150\Rightarrow 5x=120\Rightarrow x=24$: \ \rest{24 €, 48 € y 78 €}

\ex{3} $12x+8(20-x)=9{,}5\cdot20\Rightarrow 4x+160=190\Rightarrow x=7{,}5$: \ \rest{7,5 kg de 12 € y 12,5 kg de 8 €}

\ex{3} $20t+25t=90\Rightarrow\res{t=2\text{ h}}$; \ a $20\cdot2=\res{40\text{ km}}$


% =====================================================
\seccio{6}{Ecuaciones de segundo grado incompletas}

\ex{1} \begin{tasks}(1)
  \task $x^{2}-5x+6=0$: \rest{$a=1,\ b=-5,\ c=6$, completa}
  \task $3x^{2}-12=0$: \rest{$a=3,\ b=0,\ c=-12$, incompleta}
  \task $x^{2}-4x=0$: \rest{$a=1,\ b=-4,\ c=0$, incompleta}
  \task $x^{2}-4=0$: \rest{$a=1,\ b=0,\ c=-4$, incompleta}
\end{tasks}

\ex{1} \begin{tasks}(2)
  \task $x^{2}=36\ \Rightarrow\ \res{x=\pm6}$
  \task $x^{2}=25\ \Rightarrow\ \res{x=\pm5}$
  \task $x^{2}=-9\ \Rightarrow\ \res{\text{sin solución}}$
  \task $x^{2}=\tfrac{9}{4}\ \Rightarrow\ \res{x=\pm\tfrac{3}{2}}$
  \task $x^{2}=9\ \Rightarrow\ \res{x=\pm3}$
  \task $x^{2}=7\ \Rightarrow\ \res{x=\pm\sqrt{7}}$
\end{tasks}

\ex{1} \begin{tasks}(1)
  \task $x\,(x-5)=0\ \Rightarrow\ \res{x=0,\ x=5}$
  \task $x^{2}+4x=0\ \Rightarrow\ x\,(x+4)=0\ \Rightarrow\ \res{x=-4,\ x=0}$
  \task $2x^{2}-7x=0\ \Rightarrow\ x\,(2x-7)=0\ \Rightarrow\ \res{x=0,\ x=\tfrac{7}{2}}$
  \task $x^{2}-6x=0\ \Rightarrow\ x\,(x-6)=0\ \Rightarrow\ \res{x=0,\ x=6}$
\end{tasks}

\ex{2} \begin{tasks}(1)
  \task $x^{2}-9=7\ \Rightarrow\ x^{2}=16\ \Rightarrow\ \res{x=\pm4}$
  \task $x^{2}-2x=-2x+4\ \Rightarrow\ x^{2}=4\ \Rightarrow\ \res{x=\pm2}$
  \task $4x^{2}-4x+1=1\ \Rightarrow\ x^{2}-x=0\ \Rightarrow\ x\,(x-1)=0\ \Rightarrow\ \res{x=0,\ x=1}$
  \task $x^{2}+4x=3x\ \Rightarrow\ x^{2}+x=0\ \Rightarrow\ x\,(x+1)=0\ \Rightarrow\ \res{x=-1,\ x=0}$
\end{tasks}

\ex{2} \begin{tasks}(1)
  \task $\text{m.c.m.}=3:\ \ x^{2}=36\ \Rightarrow\ x^{2}=36\ \Rightarrow\ \res{x=\pm6}$
  \task $\text{m.c.m.}=4:\ \ x^{2}-2x=0\ \Rightarrow\ x\,(x-2)=0\ \Rightarrow\ \res{x=0,\ x=2}$
  \task $\text{m.c.m.}=5:\ \ (x^{2}+1)=10\ \Rightarrow\ x^{2}+1=10\ \Rightarrow\ x^{2}=9\ \Rightarrow\ \res{x=\pm3}$
  \task $\text{m.c.m.}=6:\ \ 3(x^{2}-2)-2(x^{2}+1)=6\ \Rightarrow\ x^{2}-8=6\ \Rightarrow\ x^{2}=14\ \Rightarrow\ \res{x=\pm\sqrt{14}}$
\end{tasks}

\ex{3} Al dividir entre $x$ pierde la solución $x=0$. \ $5x^2-3x=0\Rightarrow x(5x-3)=0\Rightarrow\res{x=0,\ x=\tfrac35}$.

\ex{3} \begin{tasks}(3)
  \task P. ej.\ \rest{$x^2-16=0$}
  \task P. ej.\ \rest{$x^2-3x=0$}
  \task P. ej.\ \rest{$2x^2+x=0$}
\end{tasks}


% =====================================================
\seccio{7}{Ecuaciones de segundo grado completas}

\ex{1} \begin{tasks}(2)
  \task $\Delta=4\ \Rightarrow\ x=\tfrac{6\pm2}{2}\ \Rightarrow\ \res{x=2,\ x=4}$
  \task $\Delta=36\ \Rightarrow\ x=\tfrac{-2\pm6}{2}\ \Rightarrow\ \res{x=-4,\ x=2}$
  \task $\Delta=36\ \Rightarrow\ x=\tfrac{4\pm6}{2}\ \Rightarrow\ \res{x=-1,\ x=5}$
  \task $\Delta=1\ \Rightarrow\ x=\tfrac{-7\pm1}{2}\ \Rightarrow\ \res{x=-4,\ x=-3}$
\end{tasks}

\ex{1} \begin{tasks}(2)
  \task $\Delta=9\ \Rightarrow\ x=\tfrac{5\pm3}{4}\ \Rightarrow\ \res{x=\tfrac{1}{2},\ x=2}$
  \task $\Delta=25\ \Rightarrow\ x=\tfrac{-1\pm5}{6}\ \Rightarrow\ \res{x=-1,\ x=\tfrac{2}{3}}$
  \task $\Delta=25\ \Rightarrow\ x=\tfrac{1\pm5}{12}\ \Rightarrow\ \res{x=-\tfrac{1}{3},\ x=\tfrac{1}{2}}$
  \task $\Delta=64\ \Rightarrow\ x=\tfrac{-4\pm8}{8}\ \Rightarrow\ \res{x=-\tfrac{3}{2},\ x=\tfrac{1}{2}}$
\end{tasks}

\ex{2} \begin{tasks}(1)
  \task $x^{2}-x-6=0\ \Rightarrow\ \Delta=25\ \Rightarrow\ x=\tfrac{1\pm5}{2}\ \Rightarrow\ \res{x=-2,\ x=3}$
  \task $x^{2}+2x-3=0\ \Rightarrow\ \Delta=16\ \Rightarrow\ x=\tfrac{-2\pm4}{2}\ \Rightarrow\ \res{x=-3,\ x=1}$
  \task $\Delta=81\ \Rightarrow\ x=\tfrac{1\pm9}{2}\ \Rightarrow\ \res{x=-4,\ x=5}$
  \task $x^{2}-x-6=0\ \Rightarrow\ \Delta=25\ \Rightarrow\ x=\tfrac{1\pm5}{2}\ \Rightarrow\ \res{x=-2,\ x=3}$
\end{tasks}

\ex{2} \begin{tasks}(1)
  \task $x^{2}+3x-4=6\ \Rightarrow\ x^{2}+3x-10=0\ \Rightarrow\ \Delta=49\ \Rightarrow\ x=\tfrac{-3\pm7}{2}\ \Rightarrow\ \res{x=-5,\ x=2}$
  \task $x^{2}+4x+4=3x+10\ \Rightarrow\ x^{2}+x-6=0\ \Rightarrow\ \Delta=25\ \Rightarrow\ x=\tfrac{-1\pm5}{2}\ \Rightarrow\ \res{x=-3,\ x=2}$
  \task $x^{2}+3x=2x+6\ \Rightarrow\ x^{2}+x-6=0\ \Rightarrow\ \Delta=25\ \Rightarrow\ x=\tfrac{-1\pm5}{2}\ \Rightarrow\ \res{x=-3,\ x=2}$
  \task $2x^{2}-6x+1=9\ \Rightarrow\ x^{2}-3x-4=0\ \Rightarrow\ \Delta=25\ \Rightarrow\ x=\tfrac{3\pm5}{2}\ \Rightarrow\ \res{x=-1,\ x=4}$
\end{tasks}

\ex{3} \begin{tasks}(1)
  \task $\text{m.c.m.}=6:\ \ 3x^{2}-2x=8\ \Rightarrow\ 3x^{2}-2x-8=0\ \Rightarrow\ \Delta=100\ \Rightarrow\ x=\tfrac{2\pm10}{6}\ \Rightarrow\ \res{x=-\tfrac{4}{3},\ x=2}$
  \task $\text{m.c.m.}=6:\ \ 2(x^{2}-1)+3x=12\ \Rightarrow\ 2x^{2}+3x-2=12\ \Rightarrow\ 2x^{2}+3x-14=0\ \Rightarrow\ \Delta=121\ \Rightarrow\ x=\tfrac{-3\pm11}{4}\ \Rightarrow\ \res{x=-\tfrac{7}{2},\ x=2}$
\end{tasks}

\ex{3} \begin{tasks}(1)
  \task $\text{m.c.m.}=3:\ \ x(x-1)=3x-3\ \Rightarrow\ x^{2}-x=3x-3\ \Rightarrow\ x^{2}-4x+3=0\ \Rightarrow\ \Delta=4\ \Rightarrow\ x=\tfrac{4\pm2}{2}\ \Rightarrow\ \res{x=1,\ x=3}$
  \task $\text{m.c.m.}=6:\ \ 3x(x+1)-2(x-1)(x+2)=6\ \Rightarrow\ x^{2}+x+4=6\ \Rightarrow\ x^{2}+x-2=0\ \Rightarrow\ \Delta=9\ \Rightarrow\ x=\tfrac{-1\pm3}{2}\ \Rightarrow\ \res{x=-2,\ x=1}$
\end{tasks}

\ex{3} \begin{tasks}(2)
  \task $(x-2)(x+5)=0\Rightarrow\res{x^2+3x-10=0}$
  \task $(2x-1)(x-3)=0\Rightarrow\res{2x^2-7x+3=0}$
\end{tasks}


% =====================================================
\seccio{8}{El discriminante}

\ex{1} \begin{tasks}(2)
  \task $\Delta=0$: \rest{una solución}
  \task $\Delta=-23$: \rest{ninguna solución}
  \task $\Delta=5$: \rest{dos soluciones}
  \task $\Delta=16$: \rest{dos soluciones}
\end{tasks}

\ex{1} \begin{tasks}(2)
  \task $\Delta=0$: \rest{una solución}
  \task $\Delta=49$: \rest{dos soluciones}
  \task $\Delta=-16$: \rest{ninguna solución}
  \task $\Delta=36$: \rest{dos soluciones}
  \task $\Delta=0$: \rest{una solución}
  \task $\Delta=9$: \rest{dos soluciones}
\end{tasks}

\ex{2} \begin{tasks}(1)
  \task $\Delta=0\ \Rightarrow\ x=\tfrac{6}{18}\ \Rightarrow\ \res{x=\tfrac{1}{3}\ (\text{doble})}$
  \task $\Delta=-11\ \Rightarrow\ \res{\text{sin solución}}$
  \task $\Delta=12\ \Rightarrow\ x=\tfrac{2\pm\sqrt{12}}{2}\ \Rightarrow\ \res{x=1\pm\sqrt{3}}$
  \task $\Delta=12\ \Rightarrow\ x=\tfrac{-4\pm\sqrt{12}}{2}\ \Rightarrow\ \res{x=-2\pm\sqrt{3}}$
\end{tasks}

\ex{2} $\Delta=36-4c$. \ a) $c<9$, por ejemplo \rest{$c=0$} \quad b) \rest{$c=9$} \quad c) $c>9$, por ejemplo \rest{$c=10$}

\ex{3} \begin{tasks}(1)
  \task $\Delta=16-4k=0\Rightarrow\res{k=4}$; \ $x^2-4x+4=0\Rightarrow\res{x=2}$ (doble)
  \task $m^2-64=0\Rightarrow\res{m=\pm8}$ \ ($x=-4$ o $x=4$, respectivamente)
\end{tasks}

\ex{3} \begin{tasks}(1)
  \task $\Delta=4-4m>0\Rightarrow 4>4m\Rightarrow\res{m<1}$
  \task $9+3k-6=0\Rightarrow\res{k=-1}$; \ $x^2-x-6=0\Rightarrow x=3,\ \res{x=-2}$
\end{tasks}


% =====================================================
\seccio{9}{Problemas con ecuaciones de segundo grado}

\ex{1} $x^2-5=44\Rightarrow x^2=49$: \ \rest{$x=7$ o $x=-7$}

\ex{1} $x^2=144\Rightarrow\res{x=12\text{ cm}}$ ($-12$ no vale); \ $x^2=50\Rightarrow\res{x=\sqrt{50}\approx7{,}07\text{ cm}}$

\ex{1} $x(x+3)=40\Rightarrow x^2+3x-40=0\Rightarrow x=\tfrac{-3\pm13}{2}$: \ $x=5$ ($-8$ no vale): \ \rest{5 m y 8 m}

\ex{2} $x(x+1)=156\Rightarrow x^2+x-156=0\Rightarrow x=\tfrac{-1\pm25}{2}$: \ $x=12$ ($-13$ no es natural): \ \rest{12 y 13}

\ex{2} $x^2-3x-40=0\Rightarrow x=\tfrac{3\pm13}{2}$: \ \rest{$x=8$ o $x=-5$}

\ex{2} $x^2+(x+7)^2=169\Rightarrow x^2+7x-60=0\Rightarrow x=\tfrac{-7\pm17}{2}$: \ $x=5$: \ \rest{5 cm y 12 cm}

\ex{2} $x^2-3x-130=0\Rightarrow x=\tfrac{3\pm23}{2}$: \ $x=13$ ($-10$ no vale): \ \rest{13 años}

\ex{2} $(x+3)^2-x^2=51\Rightarrow 6x+9=51\Rightarrow\res{x=7\text{ cm}}$ (los $x^2$ se cancelan: ¡es de primer grado!)

\ex{3} $x^2+(x+1)^2=85\Rightarrow x^2+x-42=0\Rightarrow x=\tfrac{-1\pm13}{2}$: \ \rest{6 y 7} \ o \ \rest{$-7$ y $-6$}

\ex{3} Si un lado mide $x$, el otro mide $17-x$: \ $x(17-x)=60\Rightarrow x^2-17x+60=0\Rightarrow x=\tfrac{17\pm7}{2}$: \ \rest{12 m y 5 m}

\ex{3} $x^2+(x+2)^2=100\Rightarrow x^2+2x-48=0\Rightarrow x=\tfrac{-2\pm14}{2}$: \ $x=6$ ($-8$ no vale): \ \rest{6 cm y 8 cm}

\ex{3} $(10+2x)(6+2x)=96\Rightarrow 4x^2+32x-36=0\Rightarrow x^2+8x-9=0\Rightarrow x=\tfrac{-8\pm10}{2}$: \ $x=1$ ($-9$ no vale): \ \rest{1 m}


\end{document}
